%15 min preso!
\documentclass[xcolor=table,aspectratio=169]{beamer}
\usepackage{beamerthemesplit}
\usepackage{wrapfig}
\usetheme{SPbGU}
\usepackage{pdfpages}
\usepackage{amsmath}
\usepackage{cmap}
\usepackage[T2A]{fontenc}
\usepackage[utf8]{inputenc}
\usepackage[english]{babel}
\usepackage{indentfirst}
\usepackage{amsmath}
\usepackage{tikz}
\usepackage{multirow}
\usepackage[noend]{algpseudocode}
\usepackage{algorithm}
\usepackage{algorithmicx}
\usepackage{fancyvrb}
\usepackage{hyperref} 
\usetikzlibrary{calc}
\usetikzlibrary{shapes, backgrounds}
\usetikzlibrary{arrows,automata}
\usetikzlibrary{positioning}
\usetikzlibrary{fit}
\usetikzlibrary{shapes.callouts}
\usetikzlibrary{shapes.misc}
\usepackage{xparse}
\usepackage{fontawesome}

\usepackage{etoolbox,refcount}
\usepackage{multicol}

\usepackage{tabularx}
\newcolumntype{Y}{>{\raggedleft\arraybackslash}X}

\renewcommand{\thealgorithm}{}

\newtheorem{mytheorem}{Theorem}
\renewcommand{\thealgorithm}{}

\newcommand{\tikzmark}[1]{\tikz[overlay,remember picture] \node (#1) {};}
\def\Put(#1,#2)#3{\leavevmode\makebox(0,0){\put(#1,#2){#3}}}

\newcommand{\ltz}{$< 1$}

\tikzset{
    state/.style={
           rectangle,
           rounded corners,
           draw=black, very thick,
           minimum height=2em,
           inner sep=2pt,
           text centered,
           },
}

\tikzset{
    invisible/.style={opacity=0,text opacity=0},
    visible on/.style={alt=#1{}{invisible}},
    alt/.code args={<#1>#2#3}{%
      \alt<#1>{\pgfkeysalso{#2}}{\pgfkeysalso{#3}} % \pgfkeysalso doesn't change the path
    },
}

\tikzset{cross/.style={cross out, draw=black, minimum size=2*(#1-\pgflinewidth), inner sep=0pt, outer sep=0pt, ultra thick},
%default radius will be 1pt. 
cross/.default={1pt}}

\NewDocumentCommand{\mycallout}{r<> O{opacity=0.8,text opacity=1} m m m}{%
\tikz[remember picture, overlay]\node[align=center, fill=cyan!20, text width=#5cm,
#2,visible on=<#1>, rounded corners,
draw,rectangle callout,anchor=pointer,callout relative pointer={(290:0.5cm)}]
at (#3) {#4};
}

\NewDocumentCommand{\mycalloutR}{r<> O{opacity=0.8,text opacity=1} m m m}{%
\tikz[remember picture, overlay]\node[align=center, fill=cyan!20, text width=#5cm,
#2,visible on=<#1>, rounded corners,
draw,rectangle callout,anchor=pointer,callout relative pointer={(30:0.8cm)}]
at (#3) {#4};
}


%callout relative pointer={(230:0.5cm)}]

\newcounter{countitems}
\newcounter{nextitemizecount}
\newcommand{\setupcountitems}{%
  \stepcounter{nextitemizecount}%
  \setcounter{countitems}{0}%
  \preto\item{\stepcounter{countitems}}%
}
\makeatletter
\newcommand{\computecountitems}{%
  \edef\@currentlabel{\number\c@countitems}%
  \label{countitems@\number\numexpr\value{nextitemizecount}-1\relax}%
}
\newcommand{\nextitemizecount}{%
  \getrefnumber{countitems@\number\c@nextitemizecount}%
}
\newcommand{\previtemizecount}{%
  \getrefnumber{countitems@\number\numexpr\value{nextitemizecount}-1\relax}%
}
\makeatother    
\newenvironment{AutoMultiColItemize}{%
\ifnumcomp{\nextitemizecount}{>}{3}{\begin{multicols}{2}}{}%
\setupcountitems\begin{itemize}}%
{\end{itemize}%
\unskip\computecountitems\ifnumcomp{\previtemizecount}{>}{3}{\end{multicols}}{}}


\beamertemplatenavigationsymbolsempty

\title[FLDDA Research Group Report]{Formal Language Driven Data Analysis Research Group Report}
\institute[SPbSU]{
Saint Petersburg State University
}

% То, что в квадратных скобках, отображается в левом нижнем углу.
\author[Semyon Grigorev]{Semyon Grigorev}

\date{September 14, 2022}


%Я предлагаю сначала в общем рассказать интересы и компетенции группы, 
%что научная группа сделала за 3 месяца, 
%потом потратить 1 страницу презентации на каждого из студентов (что сделал, почему важно для Yiming, какие планы). 
%В конце перейти к планам на год.

\begin{document}
{
\begin{frame}[fragile]
  \begin{table}
  \centering
  %\includegraphics[height=1.5cm]{pictures/SPbGU_Logo.png}
  \begin{tabularx}{\linewidth}{XcX}
    \includegraphics[height=0.9cm]{pictures/hu_logo.jpeg} \hfill
    & 
    & \hfill \includegraphics[height=1.6cm]{pictures/SPbGU_Logo.png}
  \end{tabularx}
  \end{table}
  \titlepage
\end{frame}
}

\begin{frame}[fragile]
  \frametitle{Research Area}  
  \def\firstcircle{(0,0) circle (2cm)}
  \def\secondcircle{(60:2.8cm) circle (2cm)}
  \def\thirdcircle{(0:2.8cm) circle (2cm)}
  \begin{minipage}{0.2\textwidth}
    \onslide<3->{
    Applications
    \begin{itemize}
      \item Code analysis      
      \item Code querying
      \item Code parsing
    \end{itemize}}
  \end{minipage}
  \begin{minipage}{0.48\textwidth}
  \onslide<1->{
  \begin{center}  
  \begin{tikzpicture}
    \tikzmark{xPos}{}
      \begin{scope}[shift={(3cm,-5cm)}, fill opacity=0.5]
          \fill[red] \firstcircle;
          \fill[green] \secondcircle;
          \fill[blue] \thirdcircle;
          \draw \firstcircle node[left, text width=1cm] {Graph databases};
          \draw \secondcircle node [above] {Graph analysis};
          \draw \thirdcircle node [right, text width=1.8cm] {Parsing algorithms};
      \end{scope}
  \end{tikzpicture}}  
  \mycallout<2->[opacity=1]{$(xPos) + (4.5,-4.4)$}{We are here}{2.5}
\end{center}
\end{minipage}
\begin{minipage}{0.28\textwidth}
  \onslide<4->{
  Research directions
    \begin{itemize}
      \item Graph algorithms
      \begin{itemize}
        \item Dynamic graphs
        \item Linear algebra
        \item Path querying 
      \end{itemize}
      \item Formal languages
      \begin{itemize}
        \item Languages classes and properties
        \item Parsing algorithms  
        \item Formal language constrained path querying
      \end{itemize}
    \end{itemize}}
\end{minipage}
\end{frame}

\begin{frame}[fragile]
  \frametitle{Code Analysis and Querying} 
  \tikzset{dbl/.style={double,double distance=2pt}} 
  \begin{tikzpicture}[%
    >=stealth,
    node distance=5.5cm,
    on grid,
    auto
  ]
    \node (A) [text width=4cm]             
      {Huge software projects
        \begin{itemize} 
          \item Millions LOC 
          \item Complex structure
          \item Dynamic (IDE-level analysis)
        \end{itemize}};
    \onslide<2->{
    \node (B) [right of=A, text width=4.3cm] 
      {Huge graphs for analysis
      \begin{itemize}
        \item Millions of vertices
        \item Complex structure
        \item Dynamic
      \end{itemize}
      };
    }
    \onslide<3->{
    \node (C) [right of=B, text width=5cm] 
      {Graph storage
      \begin{itemize}
        \item Graph representation
        \item Query languages
        \item Query evaluation engines        
      \end{itemize}
      };
      }
      \onslide<4->{
      \node (D) [below right = 3cm and 3cm of B, text width=6cm] 
      {Graph analysis algorithms
      \begin{itemize}
        \item Performance
        \item Nontrivial techniques (esp. for dynamic graphs)
      \end{itemize}
      };
      }
      \onslide<5->{
      \node (E) [left = 6.5cm of D, text width=5cm] 
      {Linear algebra (GraphBLAS)
      \begin{itemize}
        \item Parallel (multicore CPU, GPGPU)
        \item Flexible, expressive 
      \end{itemize}
      };}
    \onslide<2->{\path[->] (A) edge (B);}
    \onslide<3->{\path[->] (B) edge (C);}
    \onslide<4->{\path[->] (B) edge (D);}
    \onslide<5->{\path[->] (D) edge (E);}
  \end{tikzpicture}
\end{frame}

\begin{frame}[fragile]
  \frametitle{Code parsing (for IDE)}  
  \tikzset{dbl/.style={double,double distance=2pt}} 
  \begin{tikzpicture}[%
    >=stealth,
    node distance=6cm,
    on grid,
    auto
  ]
    \node (A) [text width=5cm]             
      {Parsing for IDE
        \begin{itemize} 
          \item Frequent code updates
          \item Partially correct code
          \item Multiple languages support 
          \item Performance-critical
        \end{itemize}};
    \onslide<2->{
    \node (B) [right = 9cm of A, text width=4.3cm] 
      {Parsing technique
      \begin{itemize}
        \item Error recovery
        \item Reparsing
        \item Performance
        \item Flexibility
      \end{itemize}
      };
    }
    \onslide<3->{
    \node (C) [below = 4cm of A, text width=7cm] 
      {Language description
      \begin{itemize}
        \item Modern syntax support (ambiguity, formatting-sensitivity)
        \item Human-friendly
      \end{itemize}
      };
      }
      \onslide<4->{
      \node (D) [below = 4cm of B, text width=6cm] 
      {Advanced parsing algorithms
      \begin{itemize}
        \item New formal classes of languages
        \item Error recovery
        \item Incrementalization
        \item Performance        
      \end{itemize}
      };
      }      
    \onslide<2->{\path[->] (A) edge (B);}
    \onslide<3->{\path[->] (A) edge (C);}
    \onslide<4->{\path[->] (B) edge (D);
                 \path[->] (C) edge (D);
    }    
  \end{tikzpicture}   
\end{frame}

\begin{frame}[fragile]
  \frametitle{Results} 
  \begin{center}
    \begin{tabularx}{\textwidth}{>{\setlength\hsize{1.2\hsize}\setlength\linewidth{\hsize}}X>{\setlength\hsize{.7\hsize}\setlength\linewidth{\hsize}}X>{\setlength\hsize{1.1\hsize}\setlength\linewidth{\hsize}}X}      

Graph analysis for symbolic execution engine 
&
Research prototype
&
\vspace{-10pt}
\begin{itemize}
\item Graph extraction and update mechanism
\item Constrained shortest paths for dynamic graph
\end{itemize}\\
\hline 
\onslide<2->{
  \vspace{0pt}
Graph querying algorithms
&
\vspace{0pt}
Research prototype
&
\begin{itemize}
\item New algorithms
\item Complexity analysis
\item Performance analysis
\end{itemize}\\
\hline
}
\onslide<3->{
\vspace{0pt}
Sparse linear algebra library on GPGPU
&
\vspace{0pt}
Research prototype
&
\begin{itemize}
\item Operations implementation
\item Optimizations
\item Performance analysis
\end{itemize}\\
}
\end{tabularx}
\end{center} 
\end{frame}

\begin{frame}[fragile]
  \frametitle{Rustam Azimov (Graph analysis and querying)}
  \begin{minipage}[t]{0.48\textwidth}
    \begin{itemize}
      \item[\faCheck] Linear algebra based graph analysis algorithms research and development 
      \begin{itemize}
        \item Regular path querying
        \item Context-free path querying
        \item Multiple context-free path querying
      \end{itemize}      
      \item[\faCheck] Graph analysis algorithms evaluation
      \begin{itemize}
        \item Formal language constrained path querying
        \item Static code analysis cases
        \item Linear algebra based algorithms
      \end{itemize}      
      %\item[\faGears] PhD defense
    \end{itemize}
  \end{minipage}
  \pause
  \begin{minipage}[t]{0.48\textwidth}
    %\vspace{-5pt}
    \begin{itemize}
      \item[\faHourglassHalf] Code analysis specific dataset for graph analysis algorithms evaluation and comparison
      \begin{itemize}
        \item New graphs
        \item New scenarios
        \item New quey types
      \end{itemize}      
      \item[\faHourglassHalf] Graph analysis algorithms evaluation and comparison
      \begin{itemize}
        \item New cases
        \item New graphs
        \item New algorithms
      \end{itemize}      
    \end{itemize}
  \end{minipage}
\end{frame}

\begin{frame}[fragile]
  \frametitle{Ekaterina Shemetova (Graph analysis and querying)}
  \begin{minipage}[t]{0.48\textwidth}
    \begin{itemize}
      \item[\faCheck] Linear algebra based graph analysis algorithms research and development 
      \begin{itemize}
        \item Complexity analysis
        \item Specific cases of graphs and language classes        
      \end{itemize}
      \item[\faCheck] Graph analysis algorithms evaluation
      \begin{itemize}
        \item Formal language constrained path querying
        \item Static code analysis cases
        \item GLL-based algorithms
      \end{itemize}            
      %\item[\faGears] Dynamic graph analysis
    \end{itemize}
  \end{minipage}
  \pause
  \begin{minipage}[t]{0.48\textwidth}
    %\vspace{-30pt}
    \begin{itemize}
      \item[\faHourglassHalf] Dynamic graph analysis
      \begin{itemize}
        \item Specific algorithms for symbolic execution engine and code analysis
        \item Theoretical analysis
        \item Performance analysis        
      \end{itemize}      
      \item[\faHourglassHalf] Parsing algorithms development end evaluation
      \begin{itemize}
        \item Dynamic reparsing
        \item Error recovery        
      \end{itemize}      
    \end{itemize}
  \end{minipage}
\end{frame}

\begin{frame}[fragile]
  \frametitle{Vladimir Kutuev (Graph analysis and querying)}
  \begin{minipage}[t]{0.48\textwidth}
    \begin{itemize}
      \item[\faCheck] Graph analysis algorithms development and evaluation
      \begin{itemize}
        \item Linear algebra based algorithm
        \item Static code analysis cases
      \end{itemize}
      \item[\faCheck] Graph querying algorithm fo Neo4j development and evaluation
      \begin{itemize}
        \item Static code analysis cases
        \item GLL-based algorithms
      \end{itemize}
    \end{itemize}
  \end{minipage}
  \pause
  \begin{minipage}[t]{0.48\textwidth}
    \begin{itemize}
      \item \textbf{Scholarship will be paused}
    \end{itemize}
  \end{minipage}
\end{frame}

\begin{frame}[fragile]
  \frametitle{Denis Porsev (Graph analysis and querying)}
  \begin{minipage}[t]{0.48\textwidth}
    \begin{itemize}
      \item[\faCheck] New linear algebra based algorithm for multiple source regular path querying
      \begin{itemize}
        \item Development
        \item Implementation
        \item Evaluation
      \end{itemize}      
    \end{itemize}
  \end{minipage}
  \pause
  \begin{minipage}[t]{0.48\textwidth}
    \begin{itemize}
      \item[\faHourglassHalf] The algorithm improvements
      \begin{itemize}
        \item Performance analysis and improvements
        \item Evaluation and comparison
        \item Flexibility improvements
      \end{itemize}      
    \end{itemize}
  \end{minipage}
\end{frame}

\begin{frame}[fragile]
  \frametitle{Vlada Pogozhelskaya (Graph analysis and querying)}
  \begin{minipage}[t]{0.48\textwidth}
    \begin{itemize}
      \item[\faCheck] Neo4j-based graph analysis algorithm evaluation
      \begin{itemize}
        \item General cases
        \item Static code analysis cases
      \end{itemize}  
    \end{itemize}
  \end{minipage}
  \pause
  \begin{minipage}[t]{0.48\textwidth}
    \begin{itemize}
      \item \textbf{Scholarship will be paused}
    \end{itemize}
  \end{minipage}
\end{frame}

\begin{frame}[fragile]
  \frametitle{Ilya Epelbaum (Graph analysis and querying)}
  \begin{minipage}[t]{0.48\textwidth}
    \begin{itemize}
      \item[\faCheck] Multiple Context-Free Language constrained path querying algorithm in terms of linear algebra
      \begin{itemize}
        \item Development
        \item Implementation
        \item Evaluation
      \end{itemize}
    \end{itemize}
  \end{minipage}
  \pause
  \begin{minipage}[t]{0.48\textwidth}
    \begin{itemize}
      \item \textbf{Collaboration stopped}
    \end{itemize}
    
  \end{minipage}
\end{frame}

\begin{frame}[fragile]
  \frametitle{Egor Orachyov (Linear algebra for GPGU)}
  \begin{minipage}[t]{0.48\textwidth}
    \begin{itemize}
      \item[\faCheck] GPGPU-based sparse linear algebra library design and implementation 
      \begin{itemize}
        \item High-level design and architecture of portable extensible library 
        \item Project infrastructure
        \item High-level concepts implementation
      \end{itemize}
    \end{itemize}
  \end{minipage}
  \pause
  \begin{minipage}[t]{0.49\textwidth}
    \begin{itemize}
      \item[\faHourglassHalf] Basic low-level primitives implementation
      \begin{itemize}
        \item Data structures
        \item Operations
      \end{itemize}
      \item[\faHourglassHalf] Basic graph algorithms implementation and evaluation
      \begin{itemize}
        \item BFS
        \item TC
      \end{itemize}
      \item[\faHourglassHalf] Out-of-GPGPU-memory graphs handling
    \end{itemize}
  \end{minipage}
\end{frame}

\begin{frame}[fragile]
  \frametitle{Dmitriy Panfilenok (Linear algebra for GPGU)}
  \begin{minipage}[t]{0.48\textwidth}
    \begin{itemize}
      \item[\faCheck] Kernel fusion optimization for element-wise matrix-matrix operations
      \begin{itemize}
        \item Development
        \item Implementation
        \item Evaluation
      \end{itemize} 
    \end{itemize}
  \end{minipage}
  \pause
  \begin{minipage}[t]{0.48\textwidth}
    \begin{itemize}
      \item \textbf{Scholarship will be paused}
    \end{itemize}
  \end{minipage}
\end{frame}

\begin{frame}[fragile]
  \frametitle{Kirill Garbar (Linear algebra for GPGU)}
  \begin{minipage}[t]{0.48\textwidth}
    \begin{itemize}
      \item[\faCheck] Well-typed element-wise matrix-matrix operations: generic generalized kernels for GPPGU 
      \begin{itemize}
        \item Development
        \item Implementation
        \item Evaluation        
      \end{itemize}      
    \end{itemize}
  \end{minipage}
  \pause
  \begin{minipage}[t]{0.48\textwidth}
    \begin{itemize}
      \item[\faHourglassHalf] Applied linear algebra based graph analysis algorithms
      \begin{itemize}
        \item BFS
        \item TC
        \item Supplementary matrix-vector operations
        \item Implementation and evaluation
      \end{itemize}      
    \end{itemize}
  \end{minipage}
\end{frame}

\begin{frame}[fragile]
  \frametitle{Artem Chernikov (Linear algebra for GPGU)}
  \begin{minipage}[t]{0.48\textwidth}
    \begin{itemize}
      \item[\faCheck] GPGPU-based operations on vectors and matrices
      \begin{itemize}
        \item Collections sort
        \item Matrix transpose
      \end{itemize} 
    \end{itemize}
  \end{minipage}
  \pause
  \begin{minipage}[t]{0.48\textwidth}
    \begin{itemize}
      \item[\faHourglassHalf] Generic well-typed matrix-matrix multiplication
      \begin{itemize}
        \item Sparse matrices
        \item Evaluation and comparison
      \end{itemize} 
    \end{itemize}
  \end{minipage}
\end{frame}

\begin{frame}[fragile]
  \frametitle{Alexandra Istomina (Graph analysis and querying)}
  \begin{minipage}[t]{0.48\textwidth}
    \begin{itemize}
      \item Collaboration research in previous academic year but was paused for summer
      \item Research topic: graph analysis algorithms for special types of graphs and languages
    \end{itemize}
  \end{minipage}
  \pause
  \begin{minipage}[t]{0.48\textwidth}
    \begin{itemize}
      \item[\faHourglassHalf] Graph analysis algorithms for special types of graphs and languages
      \begin{itemize}
        \item Code analysis specific cases
      \end{itemize}
    \end{itemize}
  \end{minipage}
\end{frame}

\begin{frame}[fragile]
  \frametitle{The Plan}
  \begin{minipage}[t]{0.48\textwidth}
    Code querying for declarative code analysis
    \begin{itemize}
      \item Code querying and graph querying languages
      \begin{itemize}        
        \item CodeQL
        \item Datalog
        \item GQL
        \item \ldots
      \end{itemize}
      \item Query evaluation engines
      \begin{itemize}        
        \item Performance 
        \item Flexibility        
      \end{itemize}
      \item Graph analysis algorithms
      \begin{itemize}        
        \item Performance 
        \item Scalability
        \item Incrementalization
      \end{itemize}
    \end{itemize}
  \end{minipage}
  \pause
  \begin{minipage}[t]{0.48\textwidth}
    %\vspace{-120pt}
    Parsing techniques and algorithms
    \begin{itemize}
      \item Language specification formalisms
      %\begin{itemize}
      %  \item 
      %  \item Algorithms development and analysis
      %  \item Algorithms implementation and evaluation
      %\end{itemize}
      \item Error recovery techniques
      %\begin{itemize}
      %  \item 
      %  \item Algorithms development and analysis
      %  \item Algorithms implementation and evaluation
      %\end{itemize}
      \item Reparsing techniques
      %\begin{itemize}
      %  \item 
      %  \item Algorithms development and analysis
      %  \item Algorithms implementation and evaluation
      %\end{itemize}  
    \end{itemize}
  \end{minipage}
\end{frame}

%\begin{frame}[fragile]
%  \frametitle{Scholarships request (2022--2023 academic year, 9 months)}
%    \begin{center}
%      \begin{tabular} {l  r  r} 
%        Student & Per month & Total (9 months: September -- May)\\
%        \hline
%        Egor Orachyov & 40 000      & 360 000 \\
%        Alexandra Istomina & 40 000      & 360 000 \\
%        Kirill Garbar & 20 000     & 180 000 \\
%        Denis Porsev & 20 000      & 180 000 \\
%        \hline 
%        \textbf{Total:} & 120 000 & 1 080 000 
%      \end{tabular}
%    \end{center}
%\end{frame}
\begin{frame}[fragile]
    \frametitle{Scholarships request (2022--2023 academic year, 9 months)}
      \begin{itemize}
        \item Egor Orachyov
        \item Alexandra Istomina
        \item Kirill Garbar
        \item Denis Porsev
        \item 2-3 new students 
      \end{itemize}
  \end{frame}


\end{document}
